\documentclass[11pt]{article}
\usepackage[margin=2.5cm]{geometry}
\usepackage{amsmath,amssymb,amsthm}
\usepackage{mathtools}
\usepackage{hyperref}
\usepackage[table]{xcolor}

\newtheorem{theorem}{Theorem}
\newtheorem{lemma}[theorem]{Lemma}
\newtheorem{proposition}[theorem]{Proposition}
\newtheorem{corollary}[theorem]{Corollary}
\newtheorem{definition}[theorem]{Definition}
\newtheorem{remark}[theorem]{Remark}

\newcommand{\R}{\mathbb{R}}
\newcommand{\Q}{\mathbb{Q}}
\newcommand{\N}{\mathbb{N}}
\newcommand{\coef}{\mathsf{coeff}}
\newcommand{\supp}{\mathrm{supp}}
\newcommand{\slack}{0}

\title{PP$\to$NAP Full Proof\\
\large BD Theorem 11.2, formalized and human-checked}
\author{Zinan Huang}
\date{2026-04-24}

\begin{document}
\maketitle

\begin{abstract}
This note writes out the PP$\to$NAP construction end-to-end and
cross-references every intermediate statement with the Lean 4 formalization
in \texttt{Ripple/LPP/\{Syntactic,NAP\}.lean}.
The argument is in three stages: degradation folding (\S2), $\lambda$-lifting
(\S3), and cubing + per-pair split (\S4). The main theorem (\S5) packages
them into one statement. Every step compiles \texttt{lake build}-clean; the
intent of this document is to let a human verify the \emph{mathematical}
content without re-reading Lean tactics.
\end{abstract}

%---------------------------------------------------------------------
\section{Setup}

\subsection{Production protocols}
A \emph{quadratic production protocol} on $n$ species is a tensor
$c_{r,i,j} \in \Q_{\ge 0}$ for $r, i, j \in \{1, \dots, n\}$ with
\begin{equation}
    \sum_{r=1}^n c_{r,i,j} \;=\; 2 \qquad \forall i, j.
    \tag{P1}
\end{equation}
This is the Lean structure \verb|SynPPBalance n| (\texttt{Syntactic.lean:66}).
The associated balance field on the simplex $\Delta^{n-1}$ is
\begin{equation}
    \dot x_r \;=\; \underbrace{\sum_{i,j} c_{r,i,j} x_i x_j}_{f_r(x), \text{ production}}
        \;-\; \underbrace{2 x_r \sum_{k} x_k}_{\text{degradation}}.
    \tag{P2}
\end{equation}

\subsection{Balance tensor}
Let $\sigma(x) := \sum_k x_k$. On the simplex $\sigma = 1$, but we work
formally so we keep $\sigma$ symbolic. Fold the degradation into a single
quadratic tensor:
\begin{equation}
    q_{r,i,j} \;:=\; c_{r,i,j} - [i=r] - [j=r],
\end{equation}
where $[\cdot]$ denotes the Iverson bracket. Then
\begin{equation}
    \dot x_r \;=\; \sum_{i,j} q_{r,i,j} x_i x_j,
    \tag{P3}
\end{equation}
and (P1) gives the formal conservation law
\begin{equation}
    \sum_{r} q_{r,i,j} \;=\; 2 - 1 - 1 \;=\; 0.
    \tag{P4}
\end{equation}
The Lean structure \verb|QuadField n| (\texttt{Syntactic.lean:331})
axiomatizes exactly (P4): a tensor $q_{r,i,j}$ with
$\sum_r q_{r,i,j} = 0$, no sign constraint.

\subsection{Structural predicates}
On a \verb|QuadField|, we use two pointwise predicates:
\begin{itemize}
  \item \verb|NoSelfMix zero|: $q_{r,\slack,k} = q_{r,k,\slack} = 0$ for
    all $r, k$. (``The slack coordinate never appears as a factor.'')
  \item \verb|NoSelfSelf|: $q_{r,r,r} \le 0$ for all $r$.
    (``No positive self-self production.'')
\end{itemize}
Their conjunction is \verb|SlackStructured zero|.

\begin{lemma}[\texttt{coeff\_diag\_le\_two} + \texttt{toQuadField\_noSelfSelf}]
\label{lem:auto-noselfself}
For every production protocol $c$, the folded tensor $q$ satisfies
\verb|NoSelfSelf|.
\end{lemma}
\begin{proof}
From $c_{r,i,j} \ge 0$ and $\sum_s c_{s,r,r} = 2$ we get
$c_{r,r,r} \le \sum_s c_{s,r,r} = 2$, so
$q_{r,r,r} = c_{r,r,r} - 2 \le 0$.
\hfill(\texttt{Syntactic.lean:381--398})
\end{proof}

%---------------------------------------------------------------------
\section{$\lambda$-lift}

Fix a \emph{split index} $j_0 \in \{1, \dots, n\}$ and a \emph{slack
parameter} $\lambda \in (0, 1)$. Introduce a new variable $r$ (slack) at
index~$0$ and rename the original coordinates as $y_j := x_j$ for $j \ne j_0$,
$u := \lambda x_{j_0}$. Conversely $x_{j_0} = u/\lambda$. Slack is
$r := (1-\lambda) x_{j_0}$ so that $u + r = x_{j_0}$.

The Jacobian has entries $\alpha_i := \lambda$ if $i = j_0$, else $1$.
Pushing $\dot x$ through this change of variables and adjoining the slack
equation $\dot r = \dot x_{j_0} - \dot u$ yields a new \verb|QuadField|
on $n+1$ variables indexed by $\{0, 1, \dots, n\}$ (slot $0$ = slack):
\begin{equation}
    \widetilde q_{r',i',k'} \;=\;
        \begin{cases}
          0 & i' = \slack \text{ or } k' = \slack \\
          (1-\lambda) \, q_{j_0, I, K} / (\alpha_I \alpha_K) & r' = \slack,\ i' = I^+,\ k' = K^+ \\
          \alpha_R \, q_{R, I, K} / (\alpha_I \alpha_K) & r' = R^+,\ i' = I^+,\ k' = K^+
        \end{cases}
    \tag{$\Lambda$}
\end{equation}
where $I^+$ denotes \texttt{I.succ}. (Lean: \verb|QuadField.lambdaLift|,
\texttt{Syntactic.lean:485--548}.)

\begin{lemma}[\texttt{sum\_lambdaLiftCoeff}]
\label{lem:lift-conserv}
$\widetilde q$ satisfies $\sum_{r'} \widetilde q_{r', i', k'} = 0$.
\end{lemma}
\begin{proof}
For $i' = \slack$ or $k' = \slack$ every summand is zero. For
$i' = I^+, k' = K^+$:
\[
    \sum_{r'} \widetilde q_{r', I^+, K^+}
      \;=\; \tfrac{1}{\alpha_I \alpha_K}
            \Big[ (1 - \lambda) q_{j_0, I, K} + \sum_{R} \alpha_R q_{R, I, K} \Big].
\]
The inner bracket, using $\alpha_{j_0} = \lambda$ and $\alpha_R = 1$ for
$R \ne j_0$:
\[
    (1-\lambda) q_{j_0, I, K}
      + \sum_R q_{R, I, K} + (\lambda - 1) q_{j_0, I, K}
      \;=\; \sum_R q_{R, I, K} \;=\; 0
\]
by (P4). \hfill (\texttt{Syntactic.lean:498--542})
\end{proof}

\begin{lemma}[\texttt{lambdaLift\_slackStructured}]
\label{lem:lift-slack}
If $q$ satisfies \verb|NoSelfSelf| and $\lambda > 0$, then $\widetilde q$
satisfies \verb|SlackStructured 0|.
\end{lemma}
\begin{proof}\leavevmode
\begin{itemize}
  \item \verb|NoSelfMix 0|: direct from ($\Lambda$).
  \item \verb|NoSelfSelf|: the row $r' = \slack$ gives $\widetilde q_{\slack, \slack, \slack} = 0 \le 0$.
    For $r' = R^+$:
    \[
        \widetilde q_{R^+, R^+, R^+}
          = \tfrac{\alpha_R \, q_{R, R, R}}{\alpha_R^2}
          = \tfrac{q_{R, R, R}}{\alpha_R} \;\le\; 0
    \]
    because the numerator $\le 0$ (hypothesis) and $\alpha_R > 0$ (from $\lambda > 0$).
\end{itemize}
\hfill (\texttt{Syntactic.lean:568--605})
\end{proof}

\begin{remark}[why $\lambda < 1$]
Lemma~\ref{lem:lift-slack} only uses $\lambda > 0$. The bound $\lambda < 1$
is needed downstream for the \emph{semantic} interpretation:
$r = (1-\lambda) x_{j_0}$ must be non-negative on the simplex, so
$1 - \lambda > 0$. The end-to-end theorem (Theorem~\ref{thm:main}) therefore
requires $\lambda \in (0, 1)$.
\end{remark}

%---------------------------------------------------------------------
\section{r$^2$-trick and cubing}

\subsection{r$^2$-lift}
Multiply every production term by $x_\slack^2$ (slack squared). This
rescales time and preserves the dynamics on the simplex; the crucial
algebraic consequence is that every chain-rule quadratic term
$x_a x_b$ arising from a production derivative is now weighted by
$x_\slack^2 \cdot x_a x_b$. In monomial form, the exponent vector
\begin{equation}
    \mu_{\slack, a, b}(k) \;:=\; 2 [k = \slack] + [k = a] + [k = b]
    \tag{$\mu$}
\end{equation}
is the contribution of each PP monomial $c_{r,a,b} x_a x_b$ in $f_r$ after
the $r^2$-trick. (Lean: \verb|r2PPMonomial|, \texttt{NAP.lean:1872} et seq.)

\subsection{Cubing}
A \emph{cubed index} is $\alpha : \{0, 1, \dots, n\} \to \N$ with
$\sum_k \alpha(k) = 3$. The cubed variables $v_\alpha := c_\alpha x^\alpha$
(multinomial) satisfy
\[
    \dot v_\alpha \;=\; \sum_{s} \alpha(s) \cdot c_\alpha \cdot
        \dot x_s \cdot x^{\alpha - e_s}.
\]
Substituting $\dot x_s = \sum_{a,b} \widetilde q_{s,a,b} x_\slack^2 x_a x_b$
(that is, applying $r^2$-lift in the $\lambda$-lifted field), each
$(s, a, b)$ with $\widetilde q_{s,a,b} > 0$ contributes a \emph{positive}
degree-$6$ monomial
\[
    \alpha(s) \cdot \widetilde q_{s,a,b} \cdot x^{\alpha - e_s} \cdot x^{\mu_{\slack, a, b}}
    \;=\; \alpha(s) \cdot \widetilde q_{s,a,b} \cdot x^{\delta},
\]
where
\begin{equation}
    \delta(k) \;:=\; (\alpha(k) - [k = s]) + \mu_{\slack,a,b}(k).
    \tag{$\delta$}
\end{equation}

\subsection{Non-autocatalytic split (the key theorem)}
A \emph{split} of $\delta$ is a pair $(\beta, \gamma)$ of exponent vectors
with $\beta + \gamma = \delta$. It is \emph{balanced} if
$|\beta| = |\gamma| = 3$ and \emph{non-autocatalytic} for $\alpha$ if
$\beta \ne \alpha$ and $\gamma \ne \alpha$.

\begin{theorem}[\texttt{pp\_r2\_nap\_split}, per-pair splitter]
\label{thm:per-pair}
Let $\alpha$ be a cubed index, $s, a, b$ any species with
$\alpha(s) > 0$, $a \ne \slack$, $b \ne \slack$, and suppose
\begin{equation}
    s \ne \slack \;\Longrightarrow\; a \ne s \;\lor\; b \ne s.
    \tag{$\star$}
\end{equation}
Then there exists a balanced non-autocatalytic split of
$\delta$ from ($\delta$).
\end{theorem}
\begin{proof}[Proof sketch]\leavevmode
\begin{description}
  \item[Case A: $s = \slack$.]
    Then $\mu_{\slack, a, b}$ already has two slack factors. With
    $\alpha(\slack) > 0$, set
    \[
        \beta := e_\slack + e_a + e_{a'},\qquad
        \gamma := e_\slack + e_b + e_{b'}
    \]
    where $a', b'$ are any species for which $\delta - \beta - \gamma = 0$.
    (Handled in \texttt{pp\_r2\_zero\_source\_split\_\{distinct, doubled\}}.)

  \item[Case B: $s \ne \slack$, $a \ne s$.]
    The ``foreign pair'' is $(\slack, a)$: neither coordinate equals $s$,
    so the degree-$3$ monomial $e_\slack e_a \cdot (\text{anything})$ is
    guaranteed $\ne \alpha$ whenever $\alpha(s) > 0$.
    Invoke \verb|nap_splitting_feasibility| (a flow-matching argument)
    to complete the other half. This is the substantive combinatorial
    input; it is proven in \texttt{NAP.lean:556--658}.

  \item[Case C: $s \ne \slack$, $b \ne s$.]
    Symmetric to Case B with foreign pair $(\slack, b)$.
\end{description}
The disjunction ($\star$) ensures one of B, C applies.
\hfill (\texttt{NAP.lean:1958--1983})
\end{proof}

\begin{lemma}[\texttt{quad\_r2\_nap\_split}]
\label{lem:quad-split}
On a \verb|QuadField| $F$ satisfying \verb|NoSelfSelf|, the hypothesis
($\star$) is \emph{automatically} discharged whenever
$F.\coef_{s,a,b} > 0$.
\end{lemma}
\begin{proof}
Suppose for contradiction $a = b = s$. Then
$F.\coef_{s, a, b} = F.\coef_{s,s,s} \le 0$ by \verb|NoSelfSelf|,
contradicting positivity. So at least one of $a \ne s$, $b \ne s$.
\hfill (\texttt{NAP.lean:2046--2065})
\end{proof}

%---------------------------------------------------------------------
\section{Main theorem}

\begin{theorem}[\texttt{SynPPBalance.nap\_split\_via\_lambdaLift}]
\label{thm:main}
Let $c$ be a production protocol on $n$ species ($c_{r,i,j} \ge 0$,
$\sum_r c_{r,i,j} = 2$). Let $j_0 \in \{1, \dots, n\}$, $\lambda \in (0,1)$,
and let $\widetilde q$ be the $\lambda$-lift of the folded tensor
$q_{r,i,j} = c_{r,i,j} - [i=r] - [j=r]$.

Then for every cubed index $\alpha : \{0, \dots, n\} \to \N$ with
$|\alpha| = 3$, every source $s$ with $\alpha(s) > 0$, and every pair
$(a, b)$ with $a, b \ne \slack$ and $\widetilde q_{s, a, b} > 0$:
there exists a balanced non-autocatalytic split of the chain-rule
production monomial $\delta$ defined in~($\delta$).
\end{theorem}
\begin{proof}
Compose:
\begin{enumerate}
  \item $c$ satisfies (P1) $\Rightarrow$ $q$ satisfies \verb|NoSelfSelf|
    (Lemma~\ref{lem:auto-noselfself}).
  \item $q$ satisfies \verb|NoSelfSelf| and $\lambda \in (0, 1)$
    $\Rightarrow$ $\widetilde q$ satisfies \verb|SlackStructured 0|
    (Lemma~\ref{lem:lift-slack}; only $\lambda > 0$ used for the formal
    conclusion, $\lambda < 1$ maintained for the semantic interpretation).
  \item Given $\widetilde q_{s, a, b} > 0$ with $a, b \ne \slack$ and
    $\alpha(s) > 0$: by \verb|NoSelfSelf|, the disjunction ($\star$) holds
    (Lemma~\ref{lem:quad-split}).
  \item Apply Theorem~\ref{thm:per-pair}.
\end{enumerate}
\hfill (\texttt{NAP.lean:2088--2112})
\end{proof}

%---------------------------------------------------------------------
\section{Aggregation and protocol interpretation}

Theorem~\ref{thm:main} is \emph{pointwise}: it produces a split per
$(\alpha, s, a, b)$. The protocol-level reading is:
\begin{quote}
  For every cubed $\alpha$ and every source $s$ with $\alpha(s) > 0$,
  the $r^2$-lifted chain-rule contribution to $\dot v_\alpha$,
  $\alpha(s) \sum_{a, b} \widetilde q_{s,a,b} x^\delta$, routes into a
  weighted sum of non-autocatalytic degree-$3$ interactions.
\end{quote}
The weighted-sum reconstruction identity is already formalized:
\texttt{splitRate\_reconstructs} (per-edge) and
\texttt{finite\_sum\_cubed\_direct\_pp\_with\_rate} (whole polynomial)
in \texttt{NAP.lean:1044--1055}. Negative-coefficient entries
$\widetilde q_{s, a, b} < 0$ route through the subtraction/degradation
channel, which is non-autocatalytic by construction (each degradation
edge destroys a copy, not replicates two of the same species), so the
aggregate protocol is non-autocatalytic for every cubed $\alpha$.
This is BD Theorem~11.2.

\subsection{Why naive cross-pair matching fails (forbidden edges)}
\label{sec:forbidden}

The conservation identity $\sum_{r'} \widetilde q_{r', a, b} = 0$
(Lemma~\ref{lem:lift-conserv}) invites a tempting reading as a
\emph{transfer accountability}: every positive entry
$\widetilde q_{r_1, a, b} > 0$ should be cancelled by negative entries
$\widetilde q_{r_2, a, b} < 0$ in other rows, and one might draw a
``transfer edge'' from species $r_2$ to species $r_1$ on the monomial
$x_a x_b$ with rate $|\widetilde q_{r_2, a, b}|$.

In the dual-rail picture this naive matching produces \emph{forbidden
cross-pair edges}. Concretely: the annihilation term $-k u_i v_i$
contributes a negative entry $\widetilde q_{V_i', u_i, v_i} < 0$ in the
$V_i'$ row, while a positive entry $\widetilde q_{U_j', u_i, v_i} > 0$
may sit in the $U_j'$ row of an unrelated pair $j \ne i$. The naive
recipe would pair them: ``draw an edge from $V_i'$ to $U_j'$ on the
monomial $u_i v_i$.'' Such a cross-pair edge is forbidden — it links
two different dual-rail pairs through a monomial native to neither.

Two reasons forbidden edges break the construction:

\begin{enumerate}
  \item \emph{They can carry arbitrarily large weight.} The magnitude
    of $\widetilde q_{V_i', u_i, v_i}$ is built from the full
    annihilation rate of pair $i$, which is structurally independent
    from $\widetilde q_{U_j', u_i, v_i}$ for $j \ne i$. There is no
    bound forcing these to be commensurate, so a putative cross-pair
    transfer edge would in general demand a rate the underlying NAP
    cannot realize.

  \item \emph{Naively dropping them breaks the formula condition.} If
    we refuse to draw the forbidden edge and simply discard the
    offending negative weight, the per-row reconstruction identity
    \[
       \sum_{\text{NAP edges into row } r}(\text{rate})
       \;=\; \widetilde q_{r, a, b}
    \]
    fails — the discarded weight has nowhere to go, so the aggregate
    NAP no longer reproduces the PP rate.
\end{enumerate}

The structural fix is exactly the \verb|NoSelfSelf| hypothesis
(Lemma~\ref{lem:auto-noselfself}). It restricts where negative entries
$\widetilde q_{r, a, b} < 0$ are \emph{forced} to live: on the
diagonal $(a, b) = (s, s)$ they are unconditional, but for the
off-diagonal pairs the only negative weight comes from the formal
degradation $-2 x_r \sum_k x_k$ in (P2). After the $\lambda$-lift these
degradation entries land precisely in slots that the NAP degradation
channel absorbs (each $-2 x_r x_k$ becomes a slack-tagged subtraction
edge that destroys a copy of $r$ — no cross-pair transfer required).
\verb|NoSelfSelf| therefore guarantees that every negative weight is
\emph{routable through degradation alone}, never through a forbidden
cross-pair edge.

This is why Lemma~\ref{lem:quad-split} hinges on \verb|NoSelfSelf| and
not on the weaker conservation-only hypothesis (P4). Without it, the
per-pair splitter (Theorem~\ref{thm:per-pair}) would be obliged to
manufacture exactly the forbidden cross-pair edges that the dual-rail
semantics rejects, and the formula condition would collapse.

\subsection{What the proof actually does (no transfer edges)}
\label{sec:mechanism}

Section~\ref{sec:forbidden} explains why the naive ``transfer'' picture
is unsafe. The substantive question — how does the proof get around it?
— has a clean answer:

\begin{quote}
\emph{The construction never draws an edge from a negative entry.
Every NAP edge is sourced by a single positive entry $\widetilde q_{s,a,b} > 0$,
and conservation~(P4) does the bookkeeping for the negative entries
implicitly.}
\end{quote}

Concretely, three steps:

\begin{enumerate}
  \item \emph{Generate edges only from positive entries.} For every cubed
    $\alpha$, source $s$ with $\alpha(s) > 0$, and pair $(a, b)$ with
    $a, b \ne \slack$ and $\widetilde q_{s, a, b} > 0$, the per-pair
    theorem (Theorem~\ref{thm:per-pair}) produces a balanced split
    $(\beta, \gamma)$ of $\delta$ with $\beta, \gamma \ne \alpha$. This
    split is realized as one NAP cubed-edge with rate
    $\alpha(s) \cdot \widetilde q_{s, a, b}$.

  \item \emph{\texttt{NoSelfSelf} is the precise enabler.} The per-pair
    theorem requires hypothesis ($\star$): if $s \ne \slack$ then
    $a \ne s$ or $b \ne s$. The only way ($\star$) could fail is the
    diagonal $a = b = s$ — but \verb|NoSelfSelf| forces
    $\widetilde q_{s, s, s} \le 0$, so this case has no positive entry to
    process. \verb|NoSelfSelf| precisely excludes the configuration that
    would have demanded a forbidden edge.

  \item \emph{Conservation~(P4) closes the books.} Summing the per-edge
    contributions over all positive $\widetilde q_{s, a, b}$ in a fixed
    column $(a, b)$ yields a polynomial identity (formalized as
    \texttt{splitRate\_reconstructs} together with
    \texttt{finite\_sum\_cubed\_direct\_pp\_with\_rate},
    \texttt{NAP.lean:1044--1055}) that recovers
    $\sum_{r'} \widetilde q_{r', a, b} x^\delta$ on the
    nose~--- including the rows where $\widetilde q_{r', a, b} < 0$.
    Conservation $\sum_{r'} \widetilde q_{r', a, b} = 0$ is exactly what
    makes the negative-row weights drop out as a consequence of edge
    inputs being consumed; no separate edge ever needs to be drawn for
    them.
\end{enumerate}

So the resolution to the user's worry (``does the negative weight have to
become an edge of large weight?'') is: \emph{it never has to become an
edge at all.} It appears only as a row entry on the RHS of the
polynomial identity, balanced by conservation. The largest forbidden
weight in any column is irrelevant to the construction.

\subsection{CF$'$24 worked illustration}
\label{sec:cf24}

To make this concrete, take the $\lambda$-tricked CF$'$24 system at
$\lambda = 1/2$ from \texttt{Ripple/LPP/CF24Example.lean}. Variables are
$0 = r$ (slack), $1 = u$, $2 = z_{01}$, $3 = z_{11}$:
\[
\left\{
\begin{aligned}
\dot r       &= -4u^2 + 7\,u z_{01} - 2\,u z_{11}, \\
\dot u       &= -4u^2 + 7\,u z_{01} - 2\,u z_{11}, \\
\dot z_{01}  &=  8u^2 - 16\,u z_{01} + 32\,u z_{11} - z_{01} z_{11}, \\
\dot z_{11}  &=  2\,u z_{01} - 28\,u z_{11} + z_{01} z_{11}.
\end{aligned}
\right.
\]
The folded tensor entries (using symmetric convention
$\widetilde q_{r', i, j} = \widetilde q_{r', j, i}$ so that the coefficient
of $x_i x_j$ for $i \ne j$ is $2\,\widetilde q_{r', i, j}$) for the column
$(a, b) = (u, z_{11})$ read:
\[
\bigl(\widetilde q_{r, u, z_{11}},\ \widetilde q_{u, u, z_{11}},\ \widetilde q_{z_{01}, u, z_{11}},\ \widetilde q_{z_{11}, u, z_{11}}\bigr)
\;=\; (-1,\ -1,\ +16,\ -14).
\]
Conservation: $-1 - 1 + 16 - 14 = 0$. \checkmark

\paragraph{Naive forbidden-edge view.}
The largest negative entry is $-14$ in the $z_{11}$-row. The naive
``transfer'' picture demands an edge consuming $u \cdot z_{11}$ and
producing $z_{01}$, with rate $14$. Any such direct degree-2 reaction
either consumes $z_{11}$ \emph{and} reproduces it (autocatalytic on
$z_{11}$, since the input pair already contains $z_{11}$), or doubles
the output of $z_{01}$. Both are forbidden by NAP. If we drop the edge
entirely we lose $14$ units from the $z_{01}$-row total, and
$\dot z_{01}$ no longer matches the original PP. This is exactly the
``forbidden edges of large weight'' obstruction.

\paragraph{What the proof does instead.}
Only the single positive entry $\widetilde q_{z_{01}, u, z_{11}} = +16$
generates edges. For each cubed index $\alpha$ with $\alpha(z_{01}) > 0$,
the per-pair theorem produces a split
$\beta + \gamma = \delta(k) = \bigl(\alpha(k) - [k = z_{01}]\bigr) + \mu_{\slack, u, z_{11}}(k)$
with $\beta, \gamma \ne \alpha$. Hypothesis ($\star$) is automatic
because $a = u \ne z_{01} = s$.

The aggregation identity (\texttt{splitRate\_reconstructs} +
\texttt{finite\_sum\_cubed\_direct\_pp\_with\_rate}) then guarantees
that the resulting NAP polynomial reproduces the cubed
$\dot v_\alpha$-expansion exactly, with the negative entries
$-1, -1, -14$ in rows $r, u, z_{11}$ recovered as the bookkeeping
consequence of the input $u \cdot z_{11}$ being consumed by the
$+16$-rated edge. No edge of weight~$14$ is ever drawn.

\paragraph{Where \texttt{NoSelfSelf} earns its keep.}
On this same example the diagonal column $(a, b) = (u, u)$ has
$\widetilde q_{u, u, u} = -4 \le 0$. If hypothetically the $u$-row had
been $+8$ instead of $-4$, the per-pair theorem would have to handle
$(s, a, b) = (u, u, u)$, which violates ($\star$). \verb|NoSelfSelf|
rules this configuration out structurally — the $-4$ on the diagonal
isn't an accident, it is forced by (P1):
$c_{u, u, u} - 2 \le \sum_s c_{s, u, u} - 2 = 0$.

\paragraph{Generality.}
None of this depends on CF$'$24 numerology. Conservation (P4) and
\verb|NoSelfSelf| are purely structural consequences of (P1), the
defining axiom of a production protocol. The mechanism in
\S\ref{sec:mechanism} therefore applies to every PP, with no
example-specific fudge factors hidden in the magnitudes.

%---------------------------------------------------------------------
\section{What is formally certified}

\begin{center}
\begin{tabular}{l l}
\textbf{Statement} & \textbf{Lean name} \\ \hline
$q$ conservation (P4) & \texttt{sum\_quadCoeff} \\
$q$ has \verb|NoSelfSelf| & \texttt{toQuadField\_noSelfSelf} \\
$\widetilde q$ conservation (Lemma~\ref{lem:lift-conserv}) & \texttt{sum\_lambdaLiftCoeff} \\
$\widetilde q$ is \verb|SlackStructured| (Lemma~\ref{lem:lift-slack}) & \texttt{lambdaLift\_slackStructured} \\
Per-pair split (Theorem~\ref{thm:per-pair}) & \texttt{pp\_r2\_nap\_split} \\
Main theorem (Theorem~\ref{thm:main}) & \texttt{SynPPBalance.nap\_split\_via\_lambdaLift} \\
\end{tabular}
\end{center}

All build \texttt{lake build}-clean, no \verb|sorry|, no custom axioms along
this chain. The only combinatorial content outside this note is
\verb|nap_splitting_feasibility| (\texttt{NAP.lean:556}) — the flow-matching
argument for the degree-$3$ split existence — which is itself
\verb|sorry|-free in the committed tree.

%---------------------------------------------------------------------
\section{Hole-check log}

\begin{itemize}
  \item $\lambda > 1$: semantic red line (slack goes negative); added
    $\lambda < 1$ hypothesis to Theorem~\ref{thm:main}. Underlying NAP
    statement doesn't strictly need it, but carrying it maintains the
    BD-simplex interpretation.
  \item $s = \slack$ (slack as source): originally flagged as a risk;
    resolved by observing that the $\lambda$-lift formula gives the slack
    row the same shape as any other row (with coefficient $(1-\lambda)$
    in place of $\lambda_R$), so Theorem~\ref{thm:per-pair} Case A covers it.
  \item $F$-coefficients can be negative: resolved by the observation that
    negative entries route through degradation, which is structurally
    non-autocatalytic; the split theorem only needs to handle positive
    entries.
  \item $n$-dimensional hidden obstructions (in the spirit of Vieta
    counter-examples for Newton iteration): the proof is linear in
    per-pair contribution and uniform in $n$; there is no cross-coordinate
    coupling that could introduce such obstructions.
\end{itemize}

\end{document}
